\usepackage[author={Ming Wu},date=date{},icon=Note,color=yellow]{pdfcomment}
\usepackage{graphicx}
\usepackage{svg}
\usepackage{textcomp,mathcomp,upgreek,outlines,physics}
\usepackage[version=4]{mhchem}
\usepackage{nicefrac,amsmath,mathtools,relsize}
\usepackage{algorithm,algpseudocode}
\usepackage{eqparbox,array}
\usepackage{xcolor}
% \usepackage[linesnumbered,ruled,vlined]{algorithm2e}
\usepackage{amssymb}
% \usepackage{natbib}
\usepackage[utf8]{inputenc}
\usepackage{booktabs,longtable}
\usepackage{calc}
\usepackage{bookmark}
\usepackage{hyperref}
\usepackage{authblk}%author
\usepackage{lipsum}
\usepackage{lineno}
\usepackage{filecontents}
\linenumbers
% \usepackage{siunitx}

\newcommand{\tightlist}{}
\newcommand{\micron}{$\mathrm{\(\upmu\)m}$ }

% SVG
% \svgsetup{inkscapeversion=1,inkscapepath=./figures/, inkscape=force, inkscapedpi=600, inkscapearea=page, inkscapeformat=pdf, inkscapelatex=false}
\svgsetup{inkscapeversion=1,inkscapepath=./figures/, inkscapedpi=600, inkscapearea=page, inkscapelatex=false, extractformat={pdf,eps},clean=true}

% PDF comments
\pdfcommentsetup{final}


% subsubsubsection
\newcommand{\subsubsubsection}[1]{\paragraph{#1}\mbox{}\\}
\setcounter{secnumdepth}{4}
\setcounter{tocdepth}{4}

\newcommand{\subsubsubsubsection}[1]{\paragraph{#1}\mbox{}\\}
\setcounter{secnumdepth}{5}
\setcounter{tocdepth}{5}

%% style
\usepackage[left=1cm,right=1cm,top=2cm,bottom=2cm,paperwidth=20cm,paperheight=30cm]{geometry}% Easy control of page dimensions
% paperwidth=100cm,paperheight=300cm
\usepackage{changepage}%Alter page layout
\usepackage[font=small,labelfont=bf]{caption} %make Figure in bold, and make the caption content small.

\hypersetup{
  colorlinks   = true, %Colours links instead of ugly boxes
  urlcolor     = blue, %Colour for external hyperlinks
  linkcolor    = blue, %Colour of internal links
  citecolor   = blue  %Colour of citations
}

\newcommand{\MW}[1]{\textcolor{blue}{\textbf{(MW: #1)}}}
\newcommand{\MV}[1]{\textcolor{magenta}{\textbf{(MV: #1)}}}
\newcommand{\MA}[1]{\textcolor{orange}{\textbf{(MA: #1)}}}
\newcommand{\BM}[1]{\textcolor{purple}{\textbf{(BM: #1)}}}
\newcommand{\BG}[1]{\textcolor{violet}{\textbf{(BG: #1)}}}

\providecommand{\keywords}[1]
{
  \small
  \textbf{\textit{Keywords---}} #1
}

\newcommand*{\matr}[1]{\mathbfit{#1}}
\newcommand*{\tran}{^{\mkern-1.5mu\mathsf{T}}}
\newcommand*{\conj}[1]{\overline{#1}}
\newcommand*{\hermconj}{^{\mathsf{H}}}

\newcommand*{\Register}{\textsuperscript{\tiny\textregistered}}

\newcommand*{\figref}[2][]{% \figref [<sub caption>][<figure label>]
  \hyperref[{#2}]{%
    Fig.~\ref*{#2}%
    \ifx\\#1\\%
    \else
      \,($\mathrm{#1}$)%
    \fi
  }%
}
\newcommand*{\secref}[1]{% \secref [ref]
  \hyperref[{#1}]{%
    Section~\ref*{#1}%
  }%
}

\newcommand*{\refcite}[1]{
  \hyperref[{#1}]{%
    Ref.~\cite{#1}%
  }%
}

\newcommand*{\apref}[1]{
  \hyperref[{#1}]{%
    Appendix~\ref*{#1}%
  }%
}
\newcommand*{\tableref}[1]{%
  \hyperref[{#1}]{%
    Table~\ref*{#1}%
  }%
}
\newcommand*{\equationref}[1]{%
  \hyperref[{#1}]{%
    Equation~\ref*{#1}%
  }%
}


\newcommand*{\algoref}[1]{%
  \hyperref[{#1}]{%
    Algorithm~\ref*{#1}%
  }%
}

\newcommand{\SI}[2]{${#1} \mskip3mu \mathrm{#2}$}

% Algorithms
% Comments
\algrenewcommand{\algorithmiccomment}[1]{\hfill// \eqparbox{COMMENT\thealgorithm}{#1}}
\algnewcommand{\LongComment}[1]{\hfill// \begin{minipage}[t]{\eqboxwidth{COMMENT\thealgorithm}}#1\strut\end{minipage}}
% Fun and Var
\definecolor{functionColor}{RGB}{255, 199, 6}
\definecolor{casblue}{RGB}{0, 112, 155}
% \newcommand{\var}[1]{\text{\texttt{#1}}}
\newcommand{\func}[1]{\textcolor{BlueViolet}{\texttt{#1}}}
% \newcommand{\func}[1]{\text{\textsl{#1}}}
% ForEach
\algnewcommand\algorithmicforeach{\textbf{for each}}
\algdef{S}[FOR]{ForEach}[1]{\algorithmicforeach\ #1\ \algorithmicdo}
% Sections
\newcounter{phase}[algorithm]
\newlength{\phaserulewidth}
\newcommand{\setphaserulewidth}{\setlength{\phaserulewidth}}
\newcommand{\phase}[1]{%
  \vspace{-1.25ex}
  % Top phase rule
  \Statex\leavevmode\llap{\rule{\dimexpr\labelwidth+\labelsep}{\phaserulewidth}}\rule{\linewidth}{\phaserulewidth}
  \Statex\strut\refstepcounter{phase}\textit{Phase~\thephase~--~#1}% Phase text
  % Bottom phase rule
  \vspace{-1.25ex}\Statex\leavevmode\llap{\rule{\dimexpr\labelwidth+\labelsep}{\phaserulewidth}}\rule{\linewidth}{\phaserulewidth}}
\makeatother
\setphaserulewidth{.7pt}
